\begin{textsample}{NEG Dim 4 – pi – Score -1.00 – pi/pi\_em\_42.txt}  \label{ex:f4_neg_168}
Componentes Curriculares No campo do currículo, a área de Matemática e suas Tecnologias constitui -se, na proposta da BNCC, uma das áreas de conhecimento que possui componente curricular único, apenas a Matemática.
Nessa estratégia, expõe -se o pensamento matemático como um processo em que é possível aumentar o entendimento daquilo que nos rodeia, estimulando a capacidade de pensar dos estudantes, de realizar perguntas com fundamento, contribuindo para a bagagem cultural das pessoas, além de, logicamente, ser a base para a resolução de problemas, a partir de sua utilidade prática presente no dia a dia.
Em suma, todos lidam com números, medidas, formas, operações ; todos leem e interpretam textos e gráficos, vivenciam relações de ordem e de equivalência ; todos argumentam e tiram conclusões válidas a partir de proposições verdadeiras, fazem inferências plausíveis a partir de informações parciais ou incertas.
Em outras palavras, a ninguém é permitido dispensar o conhecimento da Matemática sem abdicar de seu bem mais precioso : a consciência nas ações.
A estruturação do currículo de Matemática para o Ensino Médio com 43 habilidades prevê que os objetos do conhecimento devem ser organizados de modo a permitir ao estudante o acesso, de maneira ordenada, aos conceitos e relações que permeiam o conhecimento matemático, em busca do desenvolvimento das competências básicas para sua formação pessoal.
Esses conteúdos podem ser organizados em cinco grandes campos ou unidades temáticas, assim definidas : Álgebra/Relações, Geometria, Grandezas e Medidas, Probabilidade e Estatística, Matemática Financeira e Matemática Computacional, que são incorporados com abordagens distintas ao longo da 1ª, 2ª e 3ª séries do Ensino Médio.
Tais unidades temáticas se estabelecem como complementaridade das aprendizagens essenciais desenvolvidas no Ensino Fundamental, em especial das unidades temáticas de números, álgebra, geometria, grandezas e medidas, e probabilidade e estatística, de forma ainda mais inter-relacionada, visando possibilitar que os estudantes consolidem uma visão mais integrada da Matemática e a vejam relacionada ao seu Projeto de Vida e a aplicações em diversas situações.
Enquanto no Ensino Fundamental a proposta de Matemática relaciona as habilidades a unidades temáticas, no Ensino Médio elas se relacionam diretamente às competências específicas da área, mantendo uma ampliação das ideias centrais que foram desenvolvidas na etapa anterior, visando ampliar o Letramento Matemático, apresentar novos conhecimentos específicos, estimular processos mais elaborados de reflexão e de abstração, que deem sustentação a modos de pensar que permitam aos estudantes formular e resolver problemas em diversos contextos com mais autonomia e recursos matemáticos.
Nesta etapa, espera -se que os estudantes desenvolvam habilidades relativas aos processos de investigação, de construção de modelos e de resolução de problemas, mobilizando modos próprios de raciocinar, representar, comunicar, argumentar e, com base em discussões e validações conjuntas, aprender conceitos e desenvolver representações e procedimentos cada vez mais complexos.
Em relação aos conteúdos tradicionalmente ensinados na área de Matemática para o Ensino Médio, observa -se que na BNCC alguns tópicos, quais sejam, Matrizes e Números Complexos, Polinômios e Geometria Analítica, foram suprimidos, tendo em vista não possuírem muita aplicabilidade no Ensino Médio, e que podem, havendo necessidade, ser aprofundados no Ensino Superior.
Outro aspecto observado na proposta desta área exata na BNCC é que há uma ênfase em temáticas relacionadas à probabilidade, estatística e matemática financeira, além da inclusão de noções da matemática computacional com a introdução de conceitos para a construção de algoritmos.
Sendo assim, a aproximação entre os conteúdos escolares e o universo da cultura, a valorização das contextualizações e a busca permanente de uma instrumentação crítica para o mundo do trabalho não constituem exatamente uma novidade entre nós.
Tais princípios servem, naturalmente, de ponto de partida para a reconfiguração que agora se realiza, tendo em vista os novos passos a serem dados para o enriquecimento da prática pedagógica.
Reiteramos que este novo Currículo deve estar especialmente atento à incorporação crítica dos inúmeros recursos tecnológicos disponíveis para a representação de dados e o tratamento das informações, na busca da transformação de informação em conhecimento.
Da mesma forma, com ele é possível estabelecer articulações e conexões entre conceitos e procedimentos tanto internos à área como com outras áreas do conhecimento, colocando as situações reais como desencadeadoras de investigações por parte de alunos e professores.
Nesse âmbito, a Matemática, em parceria com a língua materna, se constitui um recurso imprescindível para uma expressão rica, uma compreensão abrangente, uma argumentação correta, um enfrentamento \textbf{assertivo} de situações-problema, uma contextualização significativa dos temas estudados.
Quando os contextos são deixados de lado, os conteúdos estudados deslocam -se sutilmente da condição de meios para a de fins das ações docentes.
O documento de Matemática para a etapa do Ensino Médio na BNCC tem uma organização por competências específicas e habilidades que se articulam com as competências gerais e, também, com as aprendizagens já previstas para o Ensino Fundamental.
No entanto, como explicitado no documento da Base : > ( … ) As competências não têm uma ordem preestabelecida.
Elas formam um todo conectado, de modo que o desenvolvimento de uma requer, em determinadas situações, a mobilização de outras. ( … ) Por sua vez, embora cada habilidade esteja associada a determinada competência, isso não significa que ela não contribua para o desenvolvimento de outras. ( p. 530 ) É importante destacar que a seleção de habilidades resultou em maior concentração de habilidades da Competência Específica 3 ( Utilizar estratégias, conceitos, definições e procedimentos matemáticos para interpretar, construir modelos e resolver problemas em diversos contextos, analisando a plausibilidade dos resultados e a adequação das soluções propostas, de modo a construir argumentação consistente. ).
O ensino para essa competência pode desenvolver a Competência 1, se a resolução de problemas em diversos contextos incluir os das Ciências da Natureza e Humanas, questões socioeconômicas ou tecnológicas, divulgados por diferentes meios.
A colaboração dessa competência com a Competência 4 se dá pela resolução de problemas com diferentes registros de representação matemáticos ( algébrico, geométrico, estatístico, computacional etc. ), na busca de solução e comunicação de resultados dos problemas. ( CE01 ) Utilizar estratégias, conceitos e procedimentos matemáticos para interpretar situações em diversos contextos, sejam atividades cotidianas, sejam fatos das Ciências da Natureza e Humanas, ou ainda questões econômicas ou tecnológicas, divulgados por diferentes meios, de modo a consolidar uma formação científica geral.
A primeira competência específica pressupõe habilidades que preparam o estudante para uma leitura crítica frente aos problemas que impactam sua vida e do seu coletivo.
Traz também conceitos e procedimentos matemáticos necessários para uma interpretação e compreensão da realidade que o aluno está inserido.
A competência 1 apresenta a Matemática como um corpo de conhecimentos a serviço de outras áreas do conhecimento e, por isso, colabora para a formação integral do estudante.
O conhecimento de estratégias, conceitos e procedimentos matemáticos, sempre levando em consideração o contexto em que a situação está inserida, estão associados ao domínio da competência.
A compreensão do que se deseja determinar de acordo com cada situação exige a combinação de vários conhecimentos de modo apropriado para que seja possível colocar esse conjunto de ideias em ação, monitorando estratégias selecionadas em cada situação e analisando sua eficiência ; e a leitura e interpretação de textos verbais, desenhos técnicos, gráficos e imagens.
É uma competência relacionada à preparação dos jovens para construir e realizar Projetos de Vida.
Vale destacar a relação dessa competência

% matched lemmas: assertivo
\end{textsample}
